\documentclass{article}
\usepackage{amsmath,hyperref}
\begin{document}
\section*{A parent fallback can invalidate independence of door scopes}
\paragraph{Source boundary.}
The selected pinned Open-RMF writer resolves a joint in the current
model first, and then in the parent. Its stored command velocity is
indexed by the resolved Entity. This fixture consists of two named
door controls and one shared \emph{or} two distinct joints, with
matched sampled positions and current sessions. Removing B's local
joint binding, while the parent mapping remains, redirects its
allowed closing request to the same Entity as A. Adding the local
binding has the opposite effect. These are modeled data fixtures
using actual resolution and writer bodies, not a claim that default
deployment or a ROS maintenance operation has this rebind interface.

\paragraph{Same reports, different work.}
For one step, A needs its aperture deficit below a declared bound.
B requests real incremental closing displacement on its target
relative to the no-B trajectory. A command notification or an
overwritten write does not earn B service. An unrelated third
actuator executes the same native close throughout.
At fully open rest in the distinct case, both orders preserve A's
zero-deficit promise and give B a positive step. With the parent
alias, OPEN then CLOSE causes positive closure at A and violates
that promise. CLOSE then OPEN makes the last command zero, so
neither violates A's promise nor performs B's work. Observing
the current leaf positions, two truthful OPEN phases, or how many
requests ran does not distinguish these conclusions.

\paragraph{Exact boundary for the one-step model.}
Let $k\in[0,16]$ be retained closing magnitude in native
acceleration increments, and write $q=\max(k-1,0)$ for the
OPEN-only first-step deficit in units of $1/512$. With a single
aliased joint, the two final deficits are
\[
  u=\min(q+1,16) \quad\hbox{(OPEN then CLOSE)},\qquad
  v=\max(\min(k+1,16)-1,0) \quad\hbox{(CLOSE then OPEN)}.
\]
The same bound $S$ permits both named obligations iff $q<u\le S$
or $q<v\le S$. With distinct joints the condition is $q\le S$, since the
separate B joint closes independently. All outputs of the recurrence
and the actual effect count are checked against native code plus the
declared one-step ideal velocity application model.

A certificate of disjoint writers is transportable only when its
consumed bindings and current cut remain valid. Re-resolution in the
same cut exposes the invalidation; it does not establish that the
Entity or the plant remains stable afterward. Numeric identity alone
is also insufficient across allocation reuse. The prior aperture
profile law assumes independently specified leaf transitions; shared
write/cache state destroys that hypothesis before any calculation
of combined aperture. A direct order scheduler with the same
resolved scope and retained velocity ties. The law and witness are
held-in source and model evidence, not field behavior or an
independently selected mission.
\end{document}
